\documentclass[11pt,a4paper]{article}
\usepackage[T1]{fontenc}
\usepackage{lmodern,amsmath,amssymb,amsthm,mathtools,booktabs,microtype,xcolor}
\usepackage[margin=26mm]{geometry}
\usepackage[colorlinks=true,linkcolor=blue!45!black,citecolor=blue!45!black,urlcolor=blue!45!black]{hyperref}
\newtheorem{theorem}{Theorem}[section]
\newtheorem{lemma}[theorem]{Lemma}
\newtheorem{proposition}[theorem]{Proposition}
\theoremstyle{remark}\newtheorem{remark}[theorem]{Remark}
\newcommand{\Q}{\mathbb Q}\newcommand{\Z}{\mathbb Z}\newcommand{\HH}{\mathcal H}
\DeclareMathOperator{\Tr}{Tr}
\setlength{\emergencystretch}{2em}
\title{A family of quartics in 301 variables}
\author{}\date{9 November 2026}
\begin{document}\maketitle
\begin{abstract}
We specify a family of integer quartics in 301 total variables by 275
quadratic equations and complete coefficient tables. Under the arithmetic
premises of the elliptic source construction, every sufficiently large
member defines the integers in the rationals. Two height detectors use
the rational trace of a sum of elliptic abscissas. The sixteen local tests
are represented by polynomial expressions with distinct valuation classes,
using fourth and eighth powers already computed for the formal logarithm.
The coefficient family is explicit, but the sufficient parameter threshold
is ineffective because the aggregate-trace height theorem has an
ineffective additive constant. No particular defining member of this
family is selected. A fully specified 351-variable quartic is a separate
result.
\end{abstract}

\section{The exact family and its scope}
All variable counts include the input $X$. Put
\[
                     A_N=80989^{2N+1}\qquad(N\ge1).
\]
The tables \nolinkurl{quadratic-family301.json} and
\nolinkurl{polynomial-family301.json} specify 275 primitive integer
quadratic rows $Q_{N,j}$ in 301 coordinates and
\[
                           P_N=\sum_{j=1}^{275}Q_{N,j}^2.
\]
Neither $A_N$ nor $N$ is a variable of $P_N$.

\begin{theorem}\label{thm:main}
Under the arithmetic premises in Section~\ref{sec:premises}, there is
an integer $N_0\ge27464$ such that, for every $N\ge N_0$ and $x\in\Q$,
\[
 x\in\Z\quad\Longleftrightarrow\quad
 \exists\boldsymbol y\in\Q^{300},\ P_N(x,\boldsymbol y)=0.
\]
For every $N\ge1$, $P_N$ has degree exactly four, content one, and
coefficient height
\[
                           H(P_N)<2^{1520}80989^{4N+2}.
\]
For $N\ge27464$ it has 4305 nonzero coordinate monomials.
The proof does not determine a numerical value of $N_0$.
\end{theorem}

An exact family specification and an effective choice of a defining
member are different assertions. The theorem supplies the first together
with eventual definition of $\Z$; its height argument does not supply
the second. The explicit 351-variable theorem in \cite{Explicit} uses
three height detectors and an effective sum-height inequality.

\section{Arithmetic premises}\label{sec:premises}
Let $K=\Q(\rho)$, $\rho^2=2$, $d=3-\rho$, and
\[
 E:y^2=x^3-d^2x,\qquad
 P=2(1+\rho,2\rho)
   =\left(\frac{57-28\rho}{8},\frac{-660+417\rho}{32}\right).
\]
The arithmetic input contract is the one precisely stated in
\cite[Section~2]{Source} and restated in \cite[Section~2]{Explicit}:
\begin{enumerate}
\item $2E(K)=\Z P$ and $h(x([n]P))\ge n^2-6$ for every nonzero
integer $n$.
\item The five-contact height bound $h(t)\le\HH M_5(t)^{19}$, with
\[
 \HH=2^{110}210^{1800}\left(\prod_{p\le76176}p\right)^{19}+128
       <2^{2093448}.
\]
Here $M_5(u/v)$ is the product of primes greater than $76176$ with
positive odd valuation on at least one of
$v,u,u-v,u-2401v,32u+343v$, in lowest terms with $v>0$, and the
excluded contact values are $0,1,2401,-343/32$.
\item The source pole obstruction outside $2,3,17$ and its simultaneous
realization for inputs in $(0,17)$ with square reduced denominator and
the specified local squareclasses and valuation bounds at $2,3,17$.
The normalized seven-coordinate equations are retained literally.
\item The simultaneous prime-pattern and conjugate reduction-order
conditions with arbitrary fixed odd avoidance modulus, finite congruences,
and real approximation, including the quantitative weighted selection
premise stated in the cited contract. In particular the slowly decreasing
separation parameter is available after all fixed data are chosen.
\end{enumerate}
No finite coefficient computation proves these inputs. The
aggregate-trace height theorem used next is a proved consequence of
Von Buhren's family height theorem, with an ineffective additive constant,
as detailed in \cite{Aggregate}. It is not an additional conjecture.

\section{The trace height and the two local detectors}
The retained square charts give root points
\[
 Q_1=[n]P,\quad Q_2=[a]P,\quad Q_3=[t]P,\quad Q_4=[b]P,
 \qquad B=\max(|n|,|a|,|t|,|b|),
\]
with nonzero integer labels. Their interpretation for all nonzero
integer labels, without a parity restriction, is proved in
\cite[Section~2]{Baseline}. Define
\[
 S_0=\frac12\Tr_{K/\Q}\left(\sum_{i=1}^4x(Q_i)\right).
\]
There are fixed constants $M_1>0,M_2\ge1$ such that
\begin{equation}\label{eq:traceheight}
                           B^2\le M_1h(S_0)+M_2.
\end{equation}
The additive constant is not effective in the cited proof. Every root
abscissa is positive at both real embeddings, so $S_0>12$.

Let $\sigma$ denote conjugation on $K/\Q$ and let $f=1/x$, extended
by zero at the identity. The auxiliary comparison points satisfy
$U_0=H$, $V_0=H-P$, and $Q_i=U_i-V_i$ as group differences on $E$.
These are the exact addition relations of the retained point charts
in \cite{Baseline}.

Use the two nonzero rational square scales $1,S_0^2$. For each
scale $a_*$ put $t_*=a_*\lambda$, $\delta=t_*-1$ and
\[
 R_* =\frac{\prod_{\theta\in\Theta}(t_*-\theta)}{a_*^5\delta^6},
 \qquad q_m=R_*^{1-2m}\delta^{-2m-2},\qquad
 \Theta=\{0,-1/2,-1,1/2399,-32/407\}.
\]
These formulas are used on their nonzero domain; the exact system
consists of their cleared quadratic equations. The case $t_*=1$
has empty contact radical and requires no reciprocal.

Define $A_i=f(U_i)$, $B_i=\sigma f(V_i)$, and
\[
 \begin{aligned}
 C_i&=L_i-L'_i\quad(1\le i\le4),& I&=L_1-\beta,\\
 D_H&=f(H),&G&=L_3^2-L_2(L_2-2)L_3L_4+(L_2-2)^2L_4^2-1.
 \end{aligned}
\]
The alternate ratios $L'_i$ already include conjugation. The nine tests
$A_i,B_i$ have order one, $C_1,I$ have order ten, and $D_H,C_2,C_3,C_4,G$
have order forty. For an ordered list of tests, use the factor
\[
 \mathcal T(D_0,\ldots,D_{s-1})
 =1+\sum_{j=0}^{s-1}q_{m(D_j)}^{2^j}D_j^{2^{j+1}},
\]
where $m(D_j)$ is its stated order. Thus different required orders can
occur within one factor.

The detector uses the product of the factors for the following six lists:
\[
 \begin{gathered}
 (D_H,C_1,B_0),\quad(C_2,I,B_2),\quad(C_3,A_2,B_3),\\
 (C_4,A_3,B_4),\quad(G,A_4),\quad(A_1,B_1).
 \end{gathered}
\]
Each of the sixteen tests occurs exactly once. The common-place valuation
argument and the literal cleared trace equations are proved in
\cite{Packing}; each contact prime supplies every required test order
at one place. The rational logarithm approximants and the scalar point
exclusions are retained.

If a nonzero integer $J$ is divisible by the $k$th powers of both
contact radicals, the contact-height theorem bounds
$h(\lambda)$ and $h(S_0^2\lambda)$ by
$(\HH+\log3)|J|^{19/k}$. Since $\lambda\ne0$,
\[
 2h(S_0)=h(S_0^2)\le h(S_0^2\lambda)+h(\lambda).
\]
Combining this with \eqref{eq:traceheight} gives
\begin{equation}\label{eq:height}
 B^2\le C|J|^{19/k}+M_2,
 \qquad C=M_1(\HH+\log3).
\end{equation}

\section{Soundness for every sufficiently large parameter}
Put $\beta=51X^2+2504$ and $c=|a-1|$. The retained scalar circuit
uses the positive contact product
\[
 r=\beta(\beta-1)(\beta-2401)(32\beta+343)(2\beta-1)
\]
and normalizer $R_s=A_Nr/(1+80989^2r)$, with shift $R_s+9792$.
The fixed-prime argument gives
\begin{equation}\label{eq:scalar}
 c\ge80989^{L_N},\quad
 L_N=\left\lfloor\frac{187(2N-1)}8\right\rfloor+1,\quad
 h(\beta)\le\HH c^{19/24}.
\end{equation}
Its integral group identities and exceptional-point exclusions are
unaffected by the local test factors.

\begin{proof}[Proof of soundness in Theorem~\ref{thm:main}]
Set
\[
 \Delta=t^2-a(a-2)tb+(a-2)^2b^2-1.
\]
The formal comparison gives divisibility of $\Delta$ by the fortieth
powers of both contact radicals. If $\Delta\ne0$, then
$|\Delta|\le14B^4$, and \eqref{eq:height} gives
\[
                     B^2\le C14^{19/40}B^{19/10}+M_2.
\]
This is impossible for sufficiently large $B$. Equation~\eqref{eq:scalar}
ensures $B\ge c-1\to\infty$ uniformly as $N\to\infty$. Hence
$\Delta=0$ for every sufficiently large $N$.

The integral Lucas descent in \cite{Source} then gives
$\log B>|a|-3\ge c-4$. Suppose $n-\beta\ne0$, and let $J$ be
its nonzero reduced numerator. The order-ten tests imply that both
tenth powers of the contact radicals divide $J$, while
$|J|\le2B\exp(h(\beta))$. For $B^2\ge2M_2$, equation~\eqref{eq:height}
therefore implies
\[
 \log B\le10\log(2C)+19\log2+19h(\beta)
          \le10\log(2C)+19\log2+19\HH c^{19/24}.
\]
The right side is $o(c)$, contradicting the Lucas lower bound for
sufficiently large $c$. Thus $n=\beta\in\Z$. Writing $X=u/v$ in
lowest terms, $v>0$, gives $v^2\mid51$ and hence $v=1$.

All thresholds depend only on the fixed constants $C,M_2,\HH$.
Since $L_N$ increases, one integer $N_0\ge27464$ works for every
$N\ge N_0$. The dependence on $M_2$ is precisely the ineffective
part of the assertion.
\end{proof}

\section{Completeness and exact verification}
Fix an integer input and a positive integer $N$. The scalar preparation
and the Lucas construction select the same integer root labels and
nonzero root points as in \cite{Baseline}. Fixed real neighborhoods,
finite preparation depths, and an avoidance modulus are chosen before
the contact primes. The modified factors introduce no new rational
function poles and retain the same formal-depth requirements.

The simultaneous completeness proposition of \cite{Packing} applies
to both scales: at every selected contact prime the cleared trace
has a unit value congruent to $1$ or $2$; the small-prime preparation
is preserved by taking powers of its already small terms; and the
real factors tend uniformly to $1$ as $\lambda\to+\infty$.
Fresh pole witnesses therefore exist under exactly the retained
arithmetic contract. Filling the rational arithmetic circuit gives
all 275 quadratic rows. Their sum of squares vanishes, proving
completeness and Theorem~\ref{thm:main}.

All required fourth and eighth powers of the $B_i$ are explicit
linear forms in existing logarithm coordinates. The five new field
powers are $C_1^4,I^4,A_2^4,A_3^4,A_4^4$, each computed by squaring
its existing square. Since every order-forty test occupies position
zero, no square or cube of $q_{40}$ is retained. The complete accounting
and both directions of the local argument are in \cite{Packing}.
The resulting system has 301 coordinates and 275 quadratic rows.

Run the following commands from \nolinkurl{computation/quartic}:
\begin{verbatim}
python build_quartic_mixed.py
python check_quartic_mixed.py
\end{verbatim}
The independent checker verifies every coefficient of the quadratic
system and the quartic. It confirms 4320 terms when powers of the
fixed parameter are recorded separately, and 4305 ordinary monomials
for $N\ge27464$. The coefficient-height bound and a retained
coefficient equal to one give the asserted height and content.
The checks do not select $N_0$ or prove the arithmetic premises.

\begin{thebibliography}{9}
\bibitem{Packing}
\emph{Simultaneous valuation conditions and explicit quartics},
9 November 2026, accompanying local theorem, simultaneous completeness
proof, exact circuit, and independent verification.
\href{../local_tests/mixed_valuation_tests_20261109_v11.pdf}{PDF}; source: \nolinkurl{mixed_valuation_tests_20261109_v11.tex}.
\bibitem{Explicit}
\emph{An explicit quartic in 351 variables}, 9 November 2026,
specified polynomial and effective conditional projection theorem.
\href{../quartic/quartic351_20261109_v11.pdf}{PDF}; source: \nolinkurl{quartic351_20261109_v11.tex}.
\bibitem{Baseline}
\emph{A family of quartics in 317 variables from aggregate trace heights},
8 November 2026, full point-chart, trace-height, and scalar arguments.
\href{quartic317_family_20261108_v9.pdf}{PDF}; source: \nolinkurl{quartic317_family_20261108_v9.tex}.
\bibitem{Source}
\emph{A forced prime condition and an explicit quartic in 392 variables},
8 November 2026, arithmetic input contract and scalar/Lucas arguments.
\href{../quartic/quartic392_20261108_v8.pdf}{PDF}; source: \nolinkurl{quartic392_20261108_v8.tex}.
\bibitem{Aggregate}
\emph{Aggregate trace heights of elliptic multiples},
8 November 2026, proof of the linear height comparison from Von Buhren's
family height theorem.
\href{../trace_heights/aggregate_trace_height_20261108_v9.pdf}{PDF}; source: \nolinkurl{aggregate_trace_height_20261108_v9.tex}.
\bibitem{Glazer}
E.~Glazer, \emph{An explicit quartic defining the integers in the rationals},
Version~13, manuscript and coefficient data, commit
\nolinkurl{5bbd74f6a479776a332cbe036809d5dddae63720},
\url{https://github.com/elliotglazer/diophantine-z-q}.
\end{thebibliography}
\end{document}
